\chapter{Requirements and System Architecture}
\section{Functional Requirements}
\begin{longtable}{p{0.12\textwidth}p{0.35\textwidth}p{0.43\textwidth}}
\caption{Functional requirements verified from source code.}\label{tab:functional-requirements}\\
\toprule
ID & Requirement & Evidence\\
\midrule
\endfirsthead
\toprule
ID & Requirement & Evidence\\
\midrule
\endhead
FR-01 & Authenticate users with access and refresh tokens & \texttt{backend/routes/auth.js}, \texttt{frontend/src/components/AuthContext.jsx}\\
FR-02 & Restrict actions by role & \texttt{backend/middleware/role.js}, \texttt{frontend/src/routes/routeConfig.jsx}\\
FR-03 & Manage users and profile avatars & \texttt{backend/routes/users.js}, \texttt{User.js}\\
FR-04 & Manage teams and subteams & \texttt{backend/routes/teams.js}, \texttt{Team.js}\\
FR-05 & Create and advance performance cycles & \texttt{backend/routes/cycles.js}, \texttt{Cycle.js}\\
FR-06 & Create, submit, approve, revise, evaluate, and lock objectives & \texttt{backend/routes/objectives.js}, \texttt{objectiveController.js}\\
FR-07 & Track KPIs embedded inside objectives & \texttt{Objective.js}, objective KPI routes\\
FR-08 & Submit and review check-ins with attachments & \texttt{CheckIn.js}, \texttt{checkInController.js}\\
FR-09 & Manage kanban-like tasks and time tracking & \texttt{Task.js}, \texttt{frontend/src/pages/TasksPage.jsx}\\
FR-10 & Generate final evaluations and HR validations & \texttt{FinalEvaluation.js}, \texttt{finalEvaluationController.js}\\
FR-11 & Record HR decisions, improvement plans, and bonus/penalty decisions & \texttt{HRDecision.js}, \texttt{ImprovementPlan.js}, \texttt{BonusPenalty.js}\\
FR-12 & Provide AI drafts and predictions & \texttt{backend/routes/ai.js}, \texttt{ai-service/app.py}\\
FR-13 & Provide analytics, reports, audit logs, and notifications & \texttt{stats.js}, \texttt{reports.js}, \texttt{auditLog.js}, \texttt{notifications.js}\\
\bottomrule
\end{longtable}

\section{Non-Functional Requirements}
\begin{table}[H]\centering\small
\caption{Non-functional requirements and implementation evidence.}
\label{tab:nonfunctional-requirements}
\begin{tabularx}{\textwidth}{p{0.25\textwidth}YY}
\toprule
Requirement & Implementation & Limitation\\
\midrule
Security & Helmet, CORS allowlist, rate limiting, XSS cleaning, Mongo sanitization, JWT, bcrypt, RBAC & No formal penetration test or CSRF protection evidence\\
Traceability & Audit log model, audit helpers, objective activity logs, final evaluation workflow history & Audit coverage may not include every route\\
Maintainability & Route-controller-model organization, services for scoring and AI, shared frontend utilities & Some legacy duplicated evaluation modules remain\\
Usability & Role-specific routes, dashboard shell, empty/loading components, charts, modals & No accessibility audit found\\
Portability & Dockerfiles, Docker Compose, Kubernetes base and overlays & AI service is not included in Compose/Kubernetes base manifests\\
Testability & Jest, Supertest, React Testing Library, validation and scoring tests & No coverage percentage stored in repository\\
\bottomrule
\end{tabularx}
\end{table}

\section{Architecture}
The frontend and backend communicate through REST-style JSON endpoints. React Router defines the protected pages, while Axios clients attach access tokens and refresh them when a request returns \texttt{401}. The backend authenticates JWT bearer tokens, loads the current user profile, checks role permissions, validates request payloads with Joi in selected routes, and stores data through Mongoose.

\begin{figure}[H]\centering
\begin{tikzpicture}[node distance=0.75cm, every node/.style={font=\scriptsize}, use/.style={draw, ellipse, fill=PerTrackLight, minimum width=2.5cm, align=center}, actor/.style={draw, rectangle, rounded corners, fill=white, minimum width=1.8cm, align=center}]
\node[actor] (admin) {Admin};
\node[actor, below=of admin] (hr) {HR};
\node[actor, below=of hr] (tl) {Team Leader};
\node[actor, below=of tl] (col) {Collaborator};
\node[use, right=2cm of admin] (users) {Manage users\\and cycles};
\node[use, right=2cm of hr] (validate) {Validate final\\evaluations};
\node[use, right=2cm of tl] (review) {Approve goals\\and reviews};
\node[use, right=2cm of col] (goals) {Create goals\\and check-ins};
\node[use, right=1.4cm of review] (ai) {AI drafts\\and predictions};
\node[use, above=of ai] (analytics) {Analytics and\\audit reports};
\draw[-Latex] (admin) -- (users);
\draw[-Latex] (admin) -- (analytics);
\draw[-Latex] (hr) -- (validate);
\draw[-Latex] (hr) -- (analytics);
\draw[-Latex] (tl) -- (review);
\draw[-Latex] (tl) -- (ai);
\draw[-Latex] (col) -- (goals);
\draw[-Latex] (col) -- (ai);
\end{tikzpicture}
\caption{Global use-case diagram. Source: author's own work based on routes and frontend route configuration.}
\label{fig:use-case}
\end{figure}

\begin{figure}[H]\centering
\begin{tikzpicture}[node distance=0.9cm, every node/.style={font=\small}, layer/.style={draw, rounded corners, fill=PerTrackLight, minimum width=11cm, minimum height=0.75cm, align=center}]
\node[layer] (ui) {Presentation layer: React pages, route guards, dashboard shell, shared components};
\node[layer, below=of ui] (client) {Client integration: Axios clients, token refresh, cached GET requests, API wrappers};
\node[layer, below=of client] (api) {API layer: Express routes, validation middleware, RBAC middleware, controllers};
\node[layer, below=of api] (domain) {Domain layer: workflow rules, objective visibility, score calculation, review context, notifications};
\node[layer, below=of domain] (data) {Data layer: Mongoose models, MongoDB collections, local uploads};
\node[layer, below=of data] (ops) {Operations layer: Docker images, Kubernetes manifests, GitHub Actions, Argo CD};
\draw[-Latex, thick] (ui) -- (client);
\draw[-Latex, thick] (client) -- (api);
\draw[-Latex, thick] (api) -- (domain);
\draw[-Latex, thick] (domain) -- (data);
\draw[-Latex, thick] (data) -- (ops);
\end{tikzpicture}
\caption{Logical architecture. Source: author's own work based on application source code.}
\label{fig:logical-architecture}
\end{figure}

\begin{figure}[H]\centering
\begin{tikzpicture}[node distance=1cm, every node/.style={font=\scriptsize}, box/.style={draw, rounded corners, fill=white, minimum width=2.4cm, minimum height=0.75cm, align=center}, group/.style={draw, rounded corners, fill=PerTrackLight, inner sep=0.25cm}]
\node[box] (routes) {Routes};
\node[box, right=of routes] (controllers) {Controllers};
\node[box, right=of controllers] (services) {Services and utils};
\node[box, right=of services] (models) {Mongoose models};
\node[box, below=of controllers] (middleware) {Auth, role, validation, audit};
\node[box, below=of services] (ai) {AI service wrappers};
\node[box, below=of models] (mongo) {MongoDB};
\draw[-Latex] (routes) -- (controllers);
\draw[-Latex] (controllers) -- (services);
\draw[-Latex] (controllers) -- (models);
\draw[-Latex] (services) -- (models);
\draw[-Latex] (models) -- (mongo);
\draw[-Latex] (middleware) -- (routes);
\draw[-Latex] (services) -- (ai);
\end{tikzpicture}
\caption{Backend component diagram. Source: author's own work based on Express route and controller structure.}
\label{fig:component-diagram}
\end{figure}

\begin{figure}[H]\centering
\resizebox{0.96\textwidth}{!}{%
\begin{tikzpicture}[every node/.style={font=\scriptsize}, comp/.style={draw, rounded corners, fill=PerTrackLight, minimum width=2.8cm, minimum height=0.8cm, align=center}, arr/.style={-Latex, thick}]
\node[comp] (route) {React route\\role and phase guard};
\node[comp, right=0.9cm of route] (axios) {Axios client\\token and refresh};
\node[comp, right=0.9cm of axios] (api) {Express route\\auth + role middleware};
\node[comp, below=0.85cm of api] (service) {Controller/service\\business rules};
\node[comp, left=0.9cm of service] (mongo) {Mongoose model\\MongoDB};
\node[comp, right=0.9cm of service] (aisvc) {AI boundary\\LLM or Flask};
\draw[arr] (route) -- node[above]{REST JSON} (axios);
\draw[arr] (axios) -- node[above]{Bearer JWT} (api);
\draw[arr] (api) -- (service);
\draw[arr] (service) -- (mongo);
\draw[arr] (service) -- (aisvc);
\draw[arr, dashed] (api.north) to[bend left=15] node[above]{401 refresh} (axios.north);
\end{tikzpicture}}
\caption{Frontend, backend, database, and AI communication flow. Source: author's own work based on API client, routes, services, and AI files.}
\label{fig:api-communication}
\end{figure}


\begin{figure}[H]\centering
\begin{tikzpicture}[node distance=0.75cm, every node/.style={font=\scriptsize}, ent/.style={draw, rounded corners, fill=PerTrackLight, minimum width=2.35cm, minimum height=0.7cm, align=center}]
\node[ent] (user) {User};
\node[ent, right=of user] (team) {Team};
\node[ent, below=of user] (cycle) {Cycle};
\node[ent, right=of cycle] (objective) {Objective\\embedded KPIs};
\node[ent, right=of objective] (task) {Task};
\node[ent, below=of objective] (checkin) {CheckIn};
\node[ent, below=of cycle] (eval) {Evaluation};
\node[ent, right=of eval] (final) {FinalEvaluation};
\node[ent, right=of final] (hr) {HRDecision};
\node[ent, below=of final] (plan) {ImprovementPlan};
\node[ent, below=of user] (audit) [yshift=-2.4cm] {AuditLog};
\draw[-Latex] (team) -- node[above]{leader/members} (user);
\draw[-Latex] (objective) -- node[above]{owner} (user);
\draw[-Latex] (objective) -- (cycle);
\draw[-Latex] (task) -- (objective);
\draw[-Latex] (checkin) -- (objective);
\draw[-Latex] (eval) -- (user);
\draw[-Latex] (eval) -- (cycle);
\draw[-Latex] (final) -- (objective);
\draw[-Latex] (hr) -- (final);
\draw[-Latex] (plan) -- (final);
\draw[-Latex] (audit) -- (user);
\end{tikzpicture}
\caption{Database model overview. Source: author's own work based on Mongoose schemas.}
\label{fig:database-model}
\end{figure}


Together, \cref{fig:use-case,fig:api-communication} show the two sides of the architecture. The use-case view explains what each actor can do, while the communication view shows how a protected frontend action becomes an authenticated backend request and, when needed, a database or AI-service interaction.

\section{Technology Choices}
React is used for component-driven UI construction \cite{reactDocs}. Vite provides the development and production build workflow \cite{viteDocs}. Express provides routing and middleware composition \cite{expressDocs}. Mongoose defines MongoDB schemas, validation, hooks, indexes, and references \cite{mongooseSchemas}. Flask is used for the Python prediction API through route decorators and JSON request handling \cite{flaskDocs}. XGBoost and scikit-learn provide the implemented model APIs \cite{xgboostDocs,sklearnRF}.

\begin{figure}[H]\centering
\resizebox{0.96\textwidth}{!}{%
\begin{tikzpicture}[every node/.style={font=\scriptsize}, layer/.style={draw, rounded corners, fill=PerTrackLight, minimum width=3.2cm, minimum height=1.0cm, align=center}, note/.style={draw, rounded corners, fill=white, minimum width=3.2cm, minimum height=0.8cm, align=center}]
\node[layer] (fe) {Frontend\\React, Vite, Router, Axios};
\node[layer, right=0.55cm of fe] (be) {Backend\\Node.js, Express, Mongoose};
\node[layer, right=0.55cm of be] (db) {Data\\MongoDB schemas and indexes};
\node[layer, below=0.55cm of fe] (ui) {Interface support\\Chart.js, Recharts, dnd-kit};
\node[layer, below=0.55cm of be] (sec) {Security\\JWT, bcrypt, Helmet, CORS};
\node[layer, below=0.55cm of db] (ai) {AI and ML\\Flask, pandas, sklearn, XGBoost};
\node[note, below=0.55cm of ui] (test) {Testing\\Jest, Supertest, RTL};
\node[note, below=0.55cm of sec] (dev) {Delivery\\Docker, Nginx, GitHub Actions};
\node[note, below=0.55cm of ai] (ops) {Deployment\\Kubernetes, Kustomize, Argo CD};
\end{tikzpicture}}
\caption{Verified technology stack grouped by project responsibility. Source: author's own work based on manifests and deployment files.}
\label{fig:technology-stack}
\end{figure}



\cref{fig:technology-stack} groups technologies by responsibility rather than listing them as isolated definitions. The trade-off is visible: the application gains separation of concerns, but it also requires careful coordination between browser state, backend authorization, MongoDB schema evolution, and a separate AI runtime.

\begin{table}[H]\centering\small
\caption{Technology-choice table based on package manifests.}
\label{tab:technology-choice}
\begin{tabularx}{\textwidth}{p{0.22\textwidth}p{0.25\textwidth}Y}
\toprule
Layer & Verified technology & Role in project\\
\midrule
Frontend & React 18, Vite, React Router, Axios & SPA, routing, API communication, role-aware pages\\
UI analytics & Chart.js, Recharts, react-chartjs-2 & Dashboards and performance visualizations\\
Backend & Node.js, Express 4, Mongoose 7 & REST API, middleware, persistence, domain rules\\
Security & JWT, bcryptjs, Helmet, CORS, rate limit, xss-clean, mongo-sanitize & Authentication and request hardening\\
AI generation & OpenAI SDK, Google Generative AI package & LLM provider abstraction for review drafts and suggestions\\
AI prediction & Flask, pandas, scikit-learn, XGBoost, joblib & Prediction API and saved model artifacts\\
DevOps & Docker, Nginx, Kubernetes, Kustomize, Argo CD, GitHub Actions & Build, deploy, GitOps and CI/CD\\
\bottomrule
\end{tabularx}
\end{table}

\section{API Architecture}
The backend mounts 26 route modules under \texttt{/api}. The main API groups are authentication, users, team members, teams, cycles, objectives, HR decisions, notifications, feed, statistics, audit logs, meetings, AI, feedback, tasks, calendar, career, performance, reports, evaluations, PDF export, check-ins, final evaluations, improvement plans, and bonus/penalty.

\section{Chapter Conclusion}
The architecture is a practical MERN-style application with a separate Python AI service. The strongest architectural feature is the separation of domain rules into services and utilities; the main architectural limitation is that the AI prediction microservice is not fully integrated into the Docker/Kubernetes deployment files.
